% Sections 9 (numerical illustrations) and 10 (further remarks and open problems).
% Labels defined here: sec:numerics, tab:numerics, sec:further, sec:tilt, conj:tilt, sec:rough, sec:nodes,
%   sec:p1-routes (= sec:nodes).

\makeatletter
\@ifundefined{conjecture}{\theoremstyle{plain}\newtheorem{conjecture}[theorem]{Conjecture}}{}
\makeatother
\providecommand{\Ath}[1]{\textup{(A$_{#1}$)}}

\section{Numerical illustrations}\label{sec:numerics}

The computations in this section are floating-point computations; they illustrate the theorems and play no role in
their proofs. They were carried out with scripts that will be made available as supplementary material (repository
URL and commit hash: to be supplied before submission).

\emph{What is shown.} Table~\ref{tab:numerics} compares the slowly varying factor
$\ell(k)=\exp\sqrt{\log(k+e)}$, which satisfies \Ath{\vartheta_p} for every $p>0$, with the pure power $\ell\equiv1$,
in the three geometries $p=\frac12$, $p=1$ and $p=2$, all with $s=2$. By Theorem~\ref{thm:main}(ii) and
Corollary~\ref{thm:pure}, the four quantities
\[
  n\Bigl(\frac{n^pb_n}{c(p)}-1\Bigr)\qquad\text{and}\qquad n\Bigl(\frac{a_n}{2b_n}-1\Bigr),
\]
for $\ell$ and for $\ell\equiv1$, stay bounded. The Bessel orders are $\alpha=1$, $0$ and $-\frac12$. The first case
lies outside the window $|\alpha|<1$ of Theorem~\ref{thm:engine} and is reached only through the shift
(Theorem~\ref{thm:shift}, from $s=\frac32$, where $\alpha=0$). The second is the case $\alpha=0$, which needs the
logarithmic form of the Bessel model. The third lies inside the window.

\emph{How the numbers were computed.} The recurrence coefficients cannot be computed from the atoms $k\le K$ alone.
The spectral weight of $P_n$ near the hard edge sits at indices $k\asymp n^{1+2/p}$, which is about $10^{16}$ for
$p=\frac12$ and $n=1600$, far beyond any truncation. Without a tail the truncated measure has a soft edge at its
smallest atom, and its coefficients obey a different law; for example, for $p=\frac12$, $s=\frac32$, the same $\ell$ and
the atoms $k\le2\cdot10^5$ alone, $n(n^pb_n/c(p)-1)$ equals $-1.81$, $-13.3$ and $-105.9$ at $n=100$, $400$ and
$1600$. This is a property of the truncation, not a counterexample.

We therefore used the atoms $k\le K$ exactly, together with a Gauss--Legendre quadrature of the continuous tail
$\int_{K+1/2}^{T}t^{-s}\ell(t)\,\delta_{t^{-p}}\diff t$ on logarithmic panels (ratio $1.02$, eight nodes per panel),
with $T=10^{48}$ for $p=\frac12$ and $T=10^{16}$ for $p=1,2$. The Jacobi matrix of the resulting discrete measure was
computed in extended (long double) precision by the updating algorithm that adds one atom at a time and restores
tridiagonal form by plane rotations. Each entry was computed for $K=10^5$ and $K=2\cdot10^5$, and for $K=2\cdot10^5$
with $T$ reduced by a factor $10^4$. The largest change of any tabulated entry under these variations is $10^{-7}$ for
$p=\frac12$, $2\cdot10^{-8}$ for $p=1$ and $2\cdot10^{-11}$ for $p=2$, so all digits shown are stable. As a further
check, the column $\ell\equiv1$ at $p=2$ agrees in every digit shown with the closed form of
Remark~\ref{rem:pure-lommel}.

\begin{table}[ht]
\centering
\small
\begin{tabular}{rrrrr}
\toprule
 & \multicolumn{2}{c}{$n\bigl(n^pb_n/c(p)-1\bigr)$} & \multicolumn{2}{c}{$n\bigl(a_n/(2b_n)-1\bigr)$}\\
\cmidrule(lr){2-3}\cmidrule(lr){4-5}
$n$ & $\ell$ & $\ell\equiv1$ & $\ell$ & $\ell\equiv1$\\
\midrule
\multicolumn{5}{l}{$p=\frac12$, $s=2$, $\alpha=1$ (reached through the shift)}\\
 100 & $-0.4039$ & $-0.5069$ & $-0.2442$ & $-0.2408$\\
 200 & $-0.4069$ & $-0.5052$ & $-0.2470$ & $-0.2453$\\
 400 & $-0.4101$ & $-0.5039$ & $-0.2484$ & $-0.2477$\\
 800 & $-0.4133$ & $-0.5028$ & $-0.2492$ & $-0.2488$\\
1600 & $-0.4163$ & $-0.5020$ & $-0.2496$ & $-0.2494$\\
\midrule
\multicolumn{5}{l}{$p=1$, $s=2$, $\alpha=0$}\\
  50 & $-0.3821$ & $-0.4934$ & $-0.4954$ & $-0.4950$\\
 100 & $-0.3893$ & $-0.4967$ & $-0.4978$ & $-0.4975$\\
 200 & $-0.3958$ & $-0.4983$ & $-0.4989$ & $-0.4987$\\
 400 & $-0.4015$ & $-0.4992$ & $-0.4994$ & $-0.4994$\\
 800 & $-0.4065$ & $-0.4996$ & $-0.4997$ & $-0.4997$\\
\midrule
\multicolumn{5}{l}{$p=2$, $s=2$, $\alpha=-\frac12$}\\
  50 & $-0.3794$ & $-0.4938$ & $-0.9756$ & $-0.9780$\\
 100 & $-0.3868$ & $-0.4969$ & $-0.9878$ & $-0.9889$\\
 200 & $-0.3936$ & $-0.4984$ & $-0.9939$ & $-0.9944$\\
 400 & $-0.3995$ & $-0.4992$ & $-0.9970$ & $-0.9972$\\
 800 & $-0.4046$ & $-0.4996$ & $-0.9985$ & $-0.9986$\\
\bottomrule
\end{tabular}
\caption{Normalised deviations from the pointwise law for $\sigma_{p,2,\ell}$, with $\ell(k)=e^{\sqrt{\log(k+e)}}$ and
for $\ell\equiv1$, computed in extended precision from the
atoms $k\le K$ and a quadrature of the continuous tail; all entries are stable under doubling of $K$ and under a change
of the tail cut-off (see the text). The limits $-(s-1)/2$ and $-p/2$ in the columns $\ell\equiv1$ are observed, not
proved. The slow drift of the first column is the effect of $\Lambda'(\log n)$
(\S\ref{sec:tilt}). Source: supplementary script \texttt{table2}.}
\label{tab:numerics}
\end{table}

\emph{What the table shows.} All entries stay bounded, as the theorems require. For the pure powers the two
normalised deviations approach $-(s-1)/2=-\frac12$ and $-p/2$ in all three geometries. For $\ell$ the deviation of
$n^pb_n$ is smaller in absolute value, and the difference from the pure power drifts slowly towards $0$, roughly like
$\frac12\Lambda'(\log n)$ with $\Lambda(\xi)=\sqrt{\log(e^\xi+e)}$; see \S\ref{sec:tilt}. The factor $\ell$ has almost
no effect on $a_n/(2b_n)$ at order $1/n$.

\section{Further remarks and open problems}\label{sec:further}

\subsection{\texorpdfstring{The $1/n$ term}{The 1/n term}}\label{sec:tilt}
Theorem~\ref{thm:main}(ii) and Corollary~\ref{thm:pure} give the rate $O(1/n)$ but not the coefficient of $1/n$. The
computations of Section~\ref{sec:numerics}, and similar ones for other $s$ (for $p=\frac12$,
$s\in\{1.5,1.75,2,2.5\}$; for $p=1$, $s\in\{2,2.5,3\}$), suggest the following. Throughout, $a_n/(2b_n)$ pairs the
diagonal entry at $P_n$ with the off-diagonal entry $b_n$ between $P_{n-1}$ and $P_n$, as in \eqref{eq:jacobi}; in
the other pairing $a_n/(2b_{n+1})$ the coefficient $-p/2$ below becomes $+p/2$, since $b_{n+1}/b_n=1-p/n+O(n^{-2})$.

\begin{conjecture}\label{conj:tilt}
\textup{(a)} For every $p>0$ and $s>1$, the pure powers satisfy
\[
  \frac{n^pb_n(\sigma_{p,s})}{c(p)}=1-\frac{s-1}{2n}+o\Bigl(\frac1n\Bigr),\qquad
  \frac{a_n(\sigma_{p,s})}{2b_n(\sigma_{p,s})}=1-\frac{p}{2n}+o\Bigl(\frac1n\Bigr).
\]
\textup{(b)} If $\ell$ satisfies \Ath{\vartheta_p} with $\ell(k)=\exp\Lambda(\log k)$, then
\[
  \frac{n^pb_n(\sigma_{p,s,\ell})}{c(p)}=1-\frac{s-1-\Lambda'(\log n)}{2n}+o\Bigl(\frac{\Lambda'(\log n)}{n}\Bigr)
  \qquad\text{whenever}\quad n\Lambda'(\log n)\to\infty.
\]
\end{conjecture}

Part (a) holds for $(p,s)=(2,2)$, by the closed form of Remark~\ref{rem:pure-lommel}. Part (b) says that, at order
$1/n$, a slowly varying factor acts like a change of the exponent $s$ by the local logarithmic derivative
$\Lambda'(\log n)$, i.e.\ like a hard-edge Bessel order that drifts with $n$. In Table~\ref{tab:numerics} the
difference between the first two columns is $0.91$ to $0.99$ times $\frac12\Lambda'(\log n)$, and this ratio
increases with $n$ in all three geometries. The $1/n$ coefficient behaves well under the shift: by the identities of
Section~\ref{sec:shift}, (a) at $s$ implies (a) at $s+p$ provided the chain-sequence parameters of
Section~\ref{sec:shift} have increments $n^pu_n-(n-1)^pu_{n-1}=o(1/n)$, a condition we can only check numerically.
Within our method, (b) amounts to computing the main term of the comparison estimate of Section~\ref{sec:hardedge}, as
Conway and Deift do for their weight \cite[(5.5)]{CD}; this would need a Bessel parametrix with a drifting order.

\subsection{Rough weights}\label{sec:rough}
The announcement in \S\ref{sec:context} shows that the pointwise law can fail for weights $k^{-s}e^{\epsilon_k}$ with
$|\epsilon_k|\le\varepsilon$, although the averaged law holds for them. Two questions remain. First, between slow
variation and arbitrary bounded oscillation there is a range of regularity conditions for which we do not know the
answer. A simple test case is $\epsilon_k=\varepsilon\sin\log k$: its increments $\epsilon_{k+1}-\epsilon_k$ tend to
$0$, but $e^{\epsilon_k}$ is not slowly varying, so Theorem~\ref{thm:main} does not apply. Second, for bounded rough
$\epsilon_k$ one may ask whether the law holds along a subsequence of density one; we do not know.

\subsection{Other nodes, other frames, larger Bessel orders}\label{sec:nodes}\label{sec:p1-routes}
Our nodes are exact powers $k^{-p}$. For nodes $x_k$ that are regularly varying in $k$ of index $-p$ the averaged
law is known \cite{Paper1}, but the pointwise law is open. Heuristically one expects $b_n/x_n\to c(p)$ under smooth
regular variation of the nodes, and failure under rough displacements of the nodes, in analogy with the weights.

For Bessel orders $\alpha\ge1$ we pass through the shift (Theorem~\ref{thm:shift}). An RH analysis of the weight itself
for $\alpha\ge1$ would require a different comparison frame (Remark~\ref{rem:alpha-ge-1}); it would be needed to compute
the $1/n$ term directly in that range.

At $p=1$ one can also work in the symmetric variable $y=\pm x^{1/2}$ of \cite{WLXZ}, where $0$ is an interior
Fisher--Hartwig point of the band; the comparison method then covers only $1<s<2$ directly, and $1<s<3$ after folding.
We do not pursue this. For pure powers and $p\ne1$ a classical quotient parametrix also works
(Remark~\ref{rem:quotient-p1}).
